Recent studies suggest that large language model (LLM) hallucinations may influence molecular property prediction, yet the nature of this effect remains poorly characterized across different hallucination types and model scales. We introduce a controlled semantic perturbation framework with three prompt-constrained perturbation operators—Structural Phantom (SP), Property Inversion (PI), and Mechanism Fabrication (MF)—alongside a residual Contextual Confabulation (CC) category. We evaluate nine conditions on two MoleculeNet benchmarks (BBBP, $N=204$; BACE, $N=152$) using multi-seed evaluation (Llama-3.1-8B) and cross-model validation (Llama-3.3-70B). A random-permutation control isolates semantic content from stylistic priming. On BBBP, prompt-constrained structural augmentation (C4a) yields the largest observed ROC-AUC increase (0.626 vs. 0.533 baseline), though shifts were not statistically significant after correction ($p > 0.05$). Post-hoc verification reveals that C4a outputs are predominantly classified as Contextual Confabulation (0% SP), indicating the effect reflects structural topic focus rather than specific hallucination types. The 70B model exhibits reversed sensitivity (C3: 0.549 vs. 0.644 baseline). On BACE, most augmentation conditions produce significant performance degradation with distributional compression toward chance level, though structural augmentation largely preserves baseline performance. The random-permutation control indicates that observed effects depend on semantic content rather than scientific register alone. These findings demonstrate that LLM-based molecular predictions exhibit differential sensitivity to structured semantic perturbations that varies by task domain and model capacity, underscoring the need for controlled evaluation in scientific reasoning pipelines.
